\subsection*{Fields}
\label{subsec:fields}

\begin{definitionbox}{Definition (Field)}
\begin{definition}[Field]
\label{def:algebraic-field}
A \textbf{field} is a set $F$ equipped with two binary operations, addition
$+$ and multiplication $\cdot$, such that:
\begin{enumerate}[label=\textbf{F\arabic*.}]
  \item $(F,+)$ is an abelian group with identity $0$;
  \item $F\setminus\{0\}$ is an abelian group under multiplication with
        identity $1$;
  \item multiplication distributes over addition.
\end{enumerate}
\end{definition}
\end{definitionbox}

\begin{remark*}[Standard quantified statement]
\[
\operatorname{Field}(F,+,\cdot,0,1)
\Longleftrightarrow
\operatorname{AbelianGroup}(F,+)
\land
\operatorname{AbelianGroup}(F\setminus\{0\},\cdot)
\land
\operatorname{LeftDistributive}(\cdot,+,F)
\land
\operatorname{RightDistributive}(\cdot,+,F).
\]
\end{remark*}

\begin{remark*}[Definition predicate reading]
A field is an additive abelian group whose nonzero elements form a
multiplicative abelian group, with multiplication distributing over addition.
\end{remark*}

\begin{remark*}[Negated quantified statement]
\[
\neg\operatorname{Field}(F,+,\cdot,0,1)
\Longleftrightarrow
\neg\operatorname{AbelianGroup}(F,+)
\lor
\neg\operatorname{AbelianGroup}(F\setminus\{0\},\cdot)
\lor
\neg\operatorname{LeftDistributive}(\cdot,+,F)
\lor
\neg\operatorname{RightDistributive}(\cdot,+,F).
\]
\end{remark*}

\begin{remark*}[Interpretation]
Fields are the algebraic setting where addition, subtraction, multiplication,
and division by nonzero elements are all available.
\end{remark*}

\begin{dependencies}
\begin{itemize}
  \item \hyperref[def:algebraic-abelian-group]{Abelian Group}
  \item \hyperref[def:algebraic-left-distributivity]{Left Distributivity}
  \item \hyperref[def:algebraic-right-distributivity]{Right Distributivity}
\end{itemize}
\end{dependencies}

\begin{proposition}[Every Field is an Integral Domain]
\label{prop:abstract-field-is-domain}
Every field is an integral domain.
\hyperref[prf:field-is-domain]{Go to proof.}
\end{proposition}

\begin{remark*}[Standard quantified statement]
\[
\operatorname{Field}(F,+,\cdot,0,1)
\Longrightarrow
\operatorname{IntegralDomain}(F,+,\cdot,0,1).
\]
\end{remark*}

\begin{remark*}[Theorem predicate reading]
The field axioms imply the integral-domain axioms.
\end{remark*}

\begin{remark*}[Negated quantified statement]
\[
\operatorname{Field}(F,+,\cdot,0,1)
\land
\neg\operatorname{IntegralDomain}(F,+,\cdot,0,1)
\Longrightarrow \bot.
\]
\end{remark*}

\begin{remark*}[Interpretation]
The existence of inverses for nonzero elements prevents nonzero zero divisors.
\end{remark*}

\begin{dependencies}
\begin{itemize}
  \item \hyperref[def:algebraic-field]{Field}
  \item \hyperref[def:algebraic-integral-domain]{Integral Domain}
  \item \hyperref[def:inverse-element]{Inverse Element}
\end{itemize}
\end{dependencies}

\begin{remark*}[Proof TODO]
Regenerate this proof as a one-proof-per-file artifact.
\end{remark*}
